\section{Purpose, claims, and the current repository boundary}\label{sec:scope}

The aim is to make abstract and continuous mathematics usable in intuitionistic
primitive recursive arithmetic without changing that arithmetic. The main
mathematical objects in this report are metric completions, spaces of uniformly
continuous maps, compact metric structures, local sections, solutions of
differential equations, and holomorphic functions. Finite encodings are their
implementation, rather than a replacement of the subject by finite combinatorics.

There are three different, compatible meanings of an embedding. First, a fixed
primitive recursive family can compute approximations, and \ipra\ can prove the
numerical error estimates. Second, an arithmetical formula can define a sort of
program codes satisfying totality and coherence conditions. Quantification over
this sort is genuine first-order quantification inside \ipra. Third, a mathematical
interpretation identifies these names with points of an ordinary metric space.
The third statement is a semantic adequacy theorem; it must not be passed off as
an internal quantification over all sets or all real numbers.

The distinction matters in both directions. Restricting the entire project to
finite tables discards useful internal abstract structure. Conversely, introducing
an arbitrary Lean type of real numbers and proving a theorem about it does not
produce an \ipra\ derivation. The framework below keeps the three meanings
separate and supplies explicit bridges between them.

\subsection{What was actually inspected}

The audit is pinned to ProveIt commit
\texttt{5f3be02388442f16487d2dc5c2529961dd783d2b}, retrieved on 20 November 2026.
The root \texttt{README.md}, \texttt{UNIFICATION-PLAN.md}, \texttt{ENDPOINTS.md},
the definitions of \texttt{PR}, \texttt{PTm}, \texttt{PRAAx}, and \texttt{Prf},
and the relevant derivation modules were read. The source inventories and the
exact paths used are recorded in the archive. This is a source audit, not a new
execution of the repository's build or its axiom audit. The repository reports a
Lean 4.32.0 toolchain; no Lean executable was available in the present environment.

The top-level description of older packets substantially understates the current
derivation library. In particular, the repository already contains object-theory
derivations for presented categories, Yoneda components, finite gluing, regular
names, coherent towers, finite-net transport, strong minimizers, contraction
estimates, and discrete Gronwall bounds. We reuse these ideas and identify their
precise boundaries rather than announce them again as new results. The canonical
target is the ordinary judgment
\[
       \Gamma\;\vdash_{\mathsf{PRAAx}}\;\varphi,
\]
implemented by \texttt{Prf PRAAx}; this follows the repository's unification
plan \cite{repo-plan,repo-pr,repo-setoid}.

\begin{center}\small
\begin{tabularx}{\linewidth}{@{}p{.22\linewidth}XX@{}}
\toprule
Topic & Existing reusable boundary & Development in this report\\
\midrule
Abstract structures & Relative setoid interpretation and categorical equations &
Three-level interpretation, relational names, completion and extension of continuous maps\\
Quantification & Finite nets and bounded syntax &
Continuous formulas with explicit error coefficients and arithmetical cross-mesh semantics\\
Local topology & Finite masks, gluing and rational covers &
Margin certificates for continuous charts, extensional gluing and transitions\\
Differential equations & Gronwall and refinement inequalities with supplied recurrences &
An actual rational Euler program, tube invariance, residual, convergence schedule and uniqueness\\
Complex analysis & Rational complex arithmetic and finite analytic estimates &
Guarded polygonal Cauchy transforms, coefficient programs and a quantified Taylor remainder\\
\bottomrule
\end{tabularx}
\end{center}

\subsection{Status of the results}

The mathematical arguments in the following sections are developed in full at
the level of the stated hypotheses. Their numerical conclusions are presented in
a form intended for direct \ipra\ derivation. The archive also contains exact
rational reference programs, executable checks of representative instances, and
Lean source outlines with explicit outstanding obligations. It does \emph{not}
contain a completed, newly kernel-checked formalization of the analytic theorems.
An abstract contract with a field requesting a proof is an interface, not a
proof of that field. This distinction is maintained in the Lean files and in
the final claim ledger.

Several broad statements are deliberately not claimed: completeness of arbitrary
metric spaces without a rate, classical compactness as a source of algorithms,
arbitrary model existence from consistency, all holomorphic functions from a
bare differentiability predicate, unrestricted cartesian closure of all computable
maps, and a uniform reflection principle for \ipra. These omissions have specific
mathematical reasons, explained alongside the positive results.

\section{The arithmetic foundation and the three levels}\label{sec:levels}

\subsection{The base theory is unchanged}

We use first-order intuitionistic natural deduction with equality. The numerical
language has symbols generated from zero, successor, projections, composition,
and primitive recursion. Every such symbol has its defining equations. The
nonlogical axioms also include successor injectivity, nonzeroness of a successor,
and induction for quantifier-free formulas. Quantifier-free formulas may use
intuitionistic propositional connectives; no induction schema for arbitrary
existential formulas is added. These are precisely the ingredients of the
repository's \texttt{PRAAx} \cite{repo-pr}.

A useful explicit coding convention is
\[
 \langle a,b\rangle=(a+b)^2+a.
\]
The pairing is injective because the intervals
$[s^2,s^2+s]$ are disjoint as $s$ varies. Both projections are obtained by search
bounded by the code, returning zero on malformed codes. A signed integer is a
pair $(u,v)$ with value $u-v$. A rational is $(u,v,d)$ with denominator $d+1$ and
value $(u-v)/(d+1)$. Reduction to lowest terms is unnecessary. Equality and order
are tested by cross multiplication. Complex rationals are pairs of rationals.
These are all numerical codes, and every arithmetic operation used below is
primitive recursive after a default value is specified on invalid input.

Lists use the empty code zero and the constructor
$\operatorname{cons}(x,s)=1+\langle x,s\rangle$. The tail of a nonempty valid
list has smaller code. Length, indexing with a default, concatenation, bounded
folds, and finite tables can therefore be defined using bounded primitive
recursion. Coding growth can be enormous. Primitive recursiveness is the
logical requirement; efficiency is a separate question, and the executable
examples use ordinary exact rational data structures.

If $R(\vec x,i)$ has a primitive recursive characteristic function, the statement
$\forall i<N\,R(\vec x,i)$ can be represented by a single Boolean fold equal
to one. Quantifier-free induction applies to that fold. This does not license
the same maneuver when $R$ is an arbitrary arithmetical formula such as
$\exists t\,\operatorname{Trace}(e,i,t,y)$ with an unbounded witness.

\subsection{Level E: executable numerical families}

At the executable level, a mathematical parameter is a tuple $p$ of numbers,
and a name is a fixed primitive recursive map $a(p,n)$ returning a rational
point at precision $n$. A modulus $\mu(p,k)$, a net $\mathcal N(p,k)$, and an
evaluator $E(p,q,k)$ are also fixed primitive recursive definitions. The word
``fixed'' concerns the program syntax; $p$ and $k$ remain arbitrary numerical
variables. Thus the result is a uniform theorem about the entire numerical
family, not a list of isolated calculations.

An endpoint typically has the form
\begin{equation}\label{eq:ep}
 \ipra\vdash\forall p,k\,
   \bigl(H(p)\longrightarrow
   V(p,k,A(p,k),C(p,k))=1\bigr),
\end{equation}
where $A,C,V$ are fixed primitive recursive functions, $C$ is a certificate
or execution table, and $H$ states the mathematical hypotheses. The output
checker is quantifier-free. The hypothesis $H$ need not be decidable: for
example, it may assert all required rational Lipschitz inequalities for a
fixed family. The proof uses the appropriate instances of $H$ when deriving
the finite Boolean invariants. Thus general mathematical hypotheses remain
visible without becoming new logical axioms.

\subsection{Level R: internally represented objects}

An arithmetical formula $D(e)$ may say that a program with code $e$ is a
total regular name. It may contain number quantifiers. A second formula
$e\simeq f$ expresses equality of the represented objects. This gives an
interpreted sort. Its quantifiers are
\[
 (\forall x\,\phi)^I=\forall e\,(D(e)\to\phi^I),\qquad
 (\exists x\,\phi)^I=\exists e\,(D(e)\land\phi^I).
\]
No new universe is postulated. The domain and equivalence are formulas over
numbers. Operations can often be implemented by primitive recursive
\emph{code constructors}: given names $e,f$, return the code of a program
which adds their approximations. The code constructor is primitive recursive
even when the programs being composed have no common primitive recursive
running-time bound. Its correctness and conditional totality must be proved.
Section~\ref{sec:relational} gives the details.

For one-sorted equation-based theories this fits the repository's
\texttt{SetoidInterpretation.\allowbreak Presentation}: its domain and equality formulas
are unrestricted, while operations are translated by fixed numerical terms.
For a relational signature, translate each atomic predicate to its own
arithmetical formula and prove invariance under the chosen equality. The
natural-deduction proof of interpretation soundness is unchanged at the
connective and quantifier steps; the atomic congruence obligations become
additional interface fields. This many-sorted relational adapter is proved on paper here,
not supplied by the existing one-sorted equational implementation. Many sorts can be kept separate or encoded by
numerical tags. Neither device supplies an unproved source axiom.

\subsection{Level S: semantic adequacy}

In the metatheory, a regular rational Cauchy name denotes a point of a complete
metric space. A proof that the algorithm's output name denotes the intended
solution connects the arithmetic theorem to the usual mathematical statement.
This is where statements such as ``the solution is differentiable'' are read
in familiar analytic language. The corresponding internal content is the
family of quantitative rational difference-quotient estimates.

Level E names are primitive recursive in their numerical parameters. Level R
names range over the codes satisfying the chosen formulas. In the standard
model they represent at most computable objects for a computable evaluator;
they do not name every real number. Level S may also formulate the same
estimates with external arbitrary names or oracle parameters. Such an
external theorem is broader semantically, but it is not thereby a
first-order \ipra\ theorem about arbitrary oracles.

\subsection{A general interpretation principle}

\begin{theorem}[Relative interpretation, with a visible boundary]\label{thm:interpretation}
Fix a finite or effectively presented first-order intuitionistic theory $T$.
For each of its sorts supply a domain formula and an equivalence formula of
\ipra, with a derivable inhabitant if the source logic requires nonempty sorts.
For each operation supply a numerical code constructor, proofs that it
preserves the domains and equivalence, and, for each relation, an extensional
arithmetical interpretation. Suppose the interpreted axioms of $T$ are
derivable in \ipra\ from a specified context $\Gamma$.
Then every fixed $T$-derivation of $\phi$ translates to an \ipra\ derivation
of $\phi^I$ from $\Gamma$ and the domain assumptions on its free variables.
\end{theorem}

\begin{proof}
Translate terms by their code constructors and atomic formulas by the supplied
formulas. Domain closure and extensionality give, by structural induction on
each fixed source term, well-definedness under substitutions. Prove by
structural induction on each fixed formula that replacing its free variables
by equivalent domain elements preserves its interpretation. At an atomic
formula this is an interface hypothesis. Conjunction, disjunction and
implication use their introduction and elimination rules. At a universal
formula, introduce an arbitrary domain element, use the induction hypothesis,
and discharge its domain premise. At an existential formula, eliminate the
witness and its domain proof and use the same witness after transport.

Now induct on the given source derivation, in the metatheory. Hypotheses are
translated hypotheses; source axioms use their supplied derivations. The
logical rules are the corresponding intuitionistic rules with the domain
guards carried through the quantifier rules. Equality substitution uses the
formula-extensionality result. Every leaf and step is therefore an actual
\ipra\ derivation. The induction here constructs a finite derivation externally;
it does not introduce a reflection principle or prove the truth of every
\ipra\ proof inside \ipra.
\end{proof}

This theorem is the correct sense in which an abstract theory can be embedded.
It does not say that all axioms of category theory, topology, or analysis
automatically satisfy the premises. A supplied completeness operation,
effective net, chart margin, or evaluation clock is part of the interpretation
data, and its necessary laws remain explicit until proved.

\subsection{Why finite calculations do not by themselves prove an analytic theorem}

Four separate obligations recur. \emph{Existence} provides an output at each
requested precision. \emph{Coherence} shows that these outputs name one
object, rather than unrelated approximate witnesses. \emph{Extensionality}
makes the result independent of the chosen name. \emph{Adequacy} relates the
name to the mathematical property, such as the integral equation of an ODE.
An endpoint that establishes only a residual bound has not yet discharged
the other three obligations.

The literature on proof mining and constructive completions motivates this
separation, but no unrestricted metatheorem from a stronger arithmetic is used
as a black box here \cite{oconnor2007,kohlenbach-oliva,avigad-feferman}.
In particular, a primitive recursive functional of higher type need not
become a first-order primitive recursive function when its functional input
is supplied only as an arbitrary program code with a totality assertion.
Every passage to first-order computation in this report specifies the
actual evaluators, clocks, or residual totality hypotheses.
